% seminario_ComoEscribir.tex
\begin{flushleft}
    {\color{gray} Presentado por: Eric Lucero Gonzalez \hfill eric.lucero@alumnos.upm.es} \\
    {\color{gray} Impartido por: Víctor Maojo \\Aula: A3202 - UPM}
    \section*{Seminario A3 Talleres de investigación: cómo escribir un artículo y hacer una presentación}
     {\color{gray} A continuación, en breves líneas y tres estilos diferentes presento una breve descripción de mi trabajo de investigación}
\end{flushleft}
\begin{center}
    \begin{tikzpicture}
        \node[aLabel, text width=0.8\textwidth, align=left] (a) at (0,0) {
            Mi investigación doctoral se centra en el diseño de una nueva metaheurística trayectorial para optimizar el problema de asignación de pistas y rutas de aproximación en el espacio aéreo terminal (TMA). Este problema es de alta complejidad combinatoria y está sujeto a múltiples restricciones operativas. La propuesta busca superar las limitaciones de los algoritmos actuales en cuanto a eficiencia y calidad de solución en un porcentaje relevante.


        };
        \node[bLabel, below=6mm of a, text width=0.8\textwidth, align=left] (b){
            Estoy estudiando cómo ayudar a los aviones a encontrar la mejor manera de llegar a la pista del aeropuerto sin chocar ni hacer fila de más. Para eso, estoy inventando un nuevo método que le dice a cada avión por dónde ir y en qué orden, como si organizara un desfile en el cielo. Con eso, quiero que los vuelos lleguen más seguros y más rápido.
        };
        \node[cLabel, below=6mm of b, text width=0.8\textwidth, align=left] (b){
            En la grisácea llanura aérea que rodea los aeropuertos —ese umbral donde los cielos se fatigan antes del descanso—, he consagrado mis días a encontrar orden en el caos. Busco, con la paciencia de un alma exiliada, un método nuevo, un pensamiento que sepa guiar los vuelos errantes hacia su destino en tierra sin colisión ni espera, sobretodo sin espera. Porque aun el aire, como los hombres, clama por una ruta más justa.
        };
    \end{tikzpicture}
\end{center}